% Edited conversion of pp. 6--8 of the source proceedings PDF.
% The portrait is preserved as a base64-encoded source asset and decoded by LuaLaTeX.
\directlua{
  dofile("scripts/base64_asset.lua")
  decode_base64_file(
    "figures/01-shevchenko/shevchenko-portrait.jpg.b64",
    "figures/01-shevchenko/shevchenko-portrait.generated.jpg"
  )
}

\ReportHeading
  {Академік Ю.~М. Шевченко --- основоположник напряму та школи термов’язкопластичності в Україні (до 100-річчя від дня народження)}
  {М.~О. Бабешко, В.~Г. Савченко}
  {Інститут механіки ім. С.~П. Тимошенка НАН України}
  {}

\begin{wrapfigure}{l}{0.34\textwidth}
  \vspace{-0.8\baselineskip}
  \centering
  \includegraphics[width=0.31\textwidth]{figures/01-shevchenko/shevchenko-portrait.generated.jpg}
  \vspace{-0.8\baselineskip}
\end{wrapfigure}

8 липня 2026 року виповнюється 100 років від дня народження видатного українського вченого-механіка Юрія Миколайовича Шевченка. Ю.~М. Шевченко народився в Києві в родині службовців Миколи Павловича та Павліни Миколаївни Шевченків. Дитячі роки та юність майбутнього академіка були нелегкими. Його батька репресували у 1938~р. Юрію Миколайовичу довелося пережити Другу світову війну, окупацію Києва та повоєнну розруху. Проте прагнення до знань допомагало долати труднощі.

У 1946~р. Ю.~М. Шевченко закінчив Київський річковий технікум і водночас вечірню середню школу. Після технікуму його призначили техніком Дніпропетровської технічної дільниці шляху Дніпровського басейнового управління, однак невдовзі він вступив до Київського державного університету імені Тараса Шевченка. Закінчивши у 1951~р. механіко-математичний факультет університету, працював учителем Володарської середньої школи Київської області, згодом викладачем Малинського лісотехнічного технікуму Житомирської області, а від 1954~р. --- асистентом кафедри теоретичної механіки Київського політехнічного інституту.

У Київському політехнічному інституті Ю.~М. Шевченко підготував кандидатську дисертацію, яку успішно захистив у 1959~р. в Інституті будівельної механіки АН УРСР (нині --- Інститут механіки ім. С.~П. Тимошенка НАН України). Від 1961 до 2016~р. він працював в Інституті механіки: спочатку старшим науковим співробітником, а від 1972~р. --- завідувачем відділу термопластичності. У 1968~р. Ю.~М. Шевченко захистив докторську дисертацію. У 1982~р. його обрали членом-кореспондентом АН УРСР, а у 1997~р. --- академіком НАН України.

Ю.~М. Шевченко отримав фундаментальні наукові результати в галузі механіки деформівного твердого тіла: механіки термов'язкопластичного деформування і руйнування елементів конструкцій за складного неізотермічного навантаження; теорії тонких оболонок і пластин; термов'язкопластичності просторових тіл. У його працях розвинено визначальні рівняння теорії термов'язкопластичності для опису деформування металів у процесах складного неізотермічного навантаження, а також експериментальні методи цієї теорії.

Під керівництвом Ю.~М. Шевченка створено сучасну експериментально обґрунтовану теорію непружного деформування ізотропних матеріалів уздовж довільних плоских траєкторій, а також траєкторій довільної кривизни й малого кручення. Праці Ю.~М. Шевченка з теоретичного та експериментального дослідження явищ запізнення скалярних і векторних властивостей матеріалів у неізотермічних процесах складного навантаження належать до світового фонду сучасної теорії термов'язкопластичності.

Виконано цикл експериментальних досліджень закономірностей непружного деформування анізотропних матеріалів і запропоновано відповідні визначальні рівняння. У працях Ю.~М. Шевченка та його учнів створено й експериментально обґрунтовано теорію термов'язкопластичного деформування ізотропних матеріалів із використанням третього інваріанта девіатора напружень. Ця теорія враховує пластичне розпушення матеріалу та залежність діаграм деформування від виду напруженого стану.

Основні результати в галузі теорії тонких оболонок і пластин полягають у розвитку методів розв'язання задач про термов'язкопружнопластичний стан неосесиметрично навантажених шаруватих оболонок змінної товщини та осесиметричних розгалужених оболонок; у чисельному аналізі впливу геометрії траєкторій деформування і третього інваріанта девіатора напружень на результати розрахунку напружено-деформованого стану оболонок. Під керівництвом Ю.~М. Шевченка створено математичну модель поведінки тонкостінних елементів трубопроводів за термопластичного деформування та розроблено методики прогнозування робочих і екстремальних станів елементів ракетної техніки за повторного термосилового навантаження.

У галузі термов'язкопластичності просторових тіл основні результати Ю.~М. Шевченка полягають у розробленні чисельних методів розв'язання плоскої задачі термопружності з урахуванням нелінійних діаграм деформування та великих деформацій; просторових задач теплопровідності й термов'язкопластичності для тіл складної форми; у чисельному дослідженні полів температур, напружень і деформацій в ізотропних, анізотропних і неоднорідних, зокрема шаруватих, тонкостінних та оболонкових елементах конструкцій за змінного термосилового і радіаційного навантаження.

Упродовж багатьох років (1973--2003) Ю.~М. Шевченко поєднував наукову роботу в Інституті механіки з науково-педагогічною діяльністю, працюючи професором кафедри опору матеріалів Київського національного університету будівництва і архітектури. Велику увагу він приділяв підготовці молодих учених. Під його керівництвом підготовлено 29 кандидатів і 10 докторів наук, які працювали в галузі механіки деформівного твердого тіла в Україні та за кордоном.

Академік Ю.~М. Шевченко був членом організаційних комітетів та учасником багатьох наукових конференцій і симпозіумів з механіки деформівного твердого тіла в Україні. Він також брав участь у роботі міжнародних конференцій в Австрії, Молдові, Польщі, Російській Федерації, США та Японії.

Наукові результати академіка Ю.~М. Шевченка опубліковано у 15 монографіях і понад 260 статтях. Основні досягнення за різними напрямами досліджень узагальнено в його оглядових статтях, опублікованих у шеститомному виданні \emph{Успехи механики} за редакцією О.~М. Гузя.

Внесок Ю.~М. Шевченка в розвиток механіки деформівного твердого тіла відзначено Державною премією УРСР у галузі науки і техніки (1980), Державною премією України в галузі науки і техніки (1993), преміями НАН України імені М.~К. Янгеля (1985) та імені О.~М. Динника (2003).

Академік Ю.~М. Шевченко пішов із життя 7 березня 2016 року.

За понад півстоліття спілкування з нашим незабутнім наставником Юрієм Миколайовичем Шевченком у нас залишилися спогади про нього як про видатну особистість, діяльність якої була спрямована на здобуття нових знань і наукових результатів. Особливу увагу він приділяв практичному застосуванню отриманих результатів та їх упровадженню. У ньому гармонійно поєднувалися принциповість і вимогливість із доброзичливим ставленням до співробітників. Світлу пам'ять про Юрія Миколайовича Шевченка ми збережемо до кінця наших днів.

\subsection*{Список монографій Ю.~М. Шевченка}
\begin{enumerate}[leftmargin=*,label=\arabic*.,itemsep=0.25em,topsep=0.4em]
\small
\item Шевченко Ю.~Н. \emph{Термопластичность при переменных нагружениях}. Київ: Наукова думка, 1970. 288~с.
\item Шевченко Ю.~Н., Бабешко М.~Е., Пискун В.~В., Савченко В.~Г. \emph{Пространственные задачи термопластичности}. Київ: Наукова думка, 1980. 262~с.
\item Шевченко Ю.~Н., Прохоренко И.~В. Теория упругопластических оболочек при неизотермических процессах нагружения. У кн.: \emph{Методы расчета оболочек}: у 5~т. Т.~3. Київ: Наукова думка, 1981. 296~с.
\item Шевченко Ю.~Н., Терехов Р.~Г. \emph{Физические уравнения термовязкопластичности}. Київ: Наукова думка, 1982. 238~с.
\item Шевченко Ю.~Н. Численные методы решения прикладных задач. У кн.: \emph{Пространственные задачи теории упругости и пластичности}: у 6~т. Т.~6. Київ: Наукова думка, 1986. 272~с.
\item Шевченко Ю.~Н., Савченко В.~Г. Термовязкопластичность. У кн.: \emph{Механика связанных полей в элементах конструкций}: у 5~т. Т.~2. Київ: Наукова думка, 1987. 264~с.
\item Шевченко Ю.~Н., Бабешко М.~Е., Терехов Р.~Г. \emph{Термовязко-упругопластические процессы сложного деформирования элементов конструкций}. Київ: Наукова думка, 1992. 328~с.
\item Шевченко Ю.~Н., Пискун В.~В., Савченко В.~Г. \emph{Решение осесимметричной пространственной задачи термопластичности на ЭЦВМ типа М-220}. Київ: Наукова думка, 1975. 108~с.
\item Шевченко Ю.~Н., Бабешко М.~Е., Прохоренко И.~В., Пискун В.~В., Савченко В.~Г. \emph{Решение осесимметричной задачи термопластичности для тонкостенных и толстостенных тел вращения на ЕС ЭВМ}. Київ: Наукова думка, 1980. 196~с.
\item Шевченко Ю.~Н., Бабешко М.~Е., Прохоренко И.~В. \emph{Методика решения осесимметричной задачи термовязкопластичности для тонких слоистых оболочек на ЕС ЭВМ}. Київ: Наукова думка, 1981. 66~с.
\item Шевченко Ю.~Н., Савченко В.~Г., Ищенко Д.~А. \emph{Метод и программа расчета на ЭВМ нестационарной теплопроводности и упругопластического напряженно-деформированного состояния элементов конструкций типа тел вращения при осесимметричных силовых и тепловых нагрузках: первая редакция}. Методические рекомендации. Київ: Институт механики АН УССР, 1988. 43~с.
\item Шевченко Ю.~Н., Савченко В.~Г., Ищенко Д.~А., Павлычко В.~М. \emph{Расчеты и испытания на прочность. Метод и пакет прикладных программ расчета на ЭВМ нестационарной теплопроводности и упругопластического напряженно-деформированного состояния элементов конструкций типа тел вращения при осесимметричных силовых и тепловых нагрузках}. Рекомендации Р~54-284-90. Москва: ВНИИНМАШ Госстандарта СССР, 1990. 56~с.
\item Григоренко Я.~М., Шевченко Ю.~Н., Василенко А.~Т. та ін. Численные методы. У кн.: \emph{Механика композитов}: у 12~т. Т.~11. Київ: А.С.К., 2002. 448~с.
\item Шевченко Ю.~М. \emph{Розв'язок плоскої задачі теорії пружності у напруженнях методом скінченних різниць. Конспект лекцій та методика виконання розрахунково-графічної роботи ``Розрахунок на міцність балки-стінки''}. Київ: КНУБА, 2003. 52~с.
\item Шевченко Ю.~М. \emph{Теорія тонкостінних стержнів}. Конспект лекцій. Київ: КНУБА, 2003. 76~с.
\end{enumerate}
